% =====================================================================
% IBFE Tutorials: motivation for the immersed boundary finite element method
% =====================================================================
\documentclass[aspectratio=169,11pt]{beamer}

% Match the visual language of the SIAM Waves talk.
\usetheme{Boadilla}
\usecolortheme{seahorse}
\useinnertheme{rectangles}
\useoutertheme{infolines}

\setbeamercolor{structure}{fg=blue!55!black}
\setbeamercolor{title}{fg=white,bg=blue!55!black}
\setbeamercolor{frametitle}{fg=blue!55!black}
\setbeamercolor{block title}{bg=blue!12,fg=blue!55!black}
\setbeamercolor{block body}{bg=blue!4}
\setbeamertemplate{navigation symbols}{}
\setbeamerfont{frametitle}{series=\bfseries,size=\large}
\setbeamerfont{title}{series=\bfseries,size=\Large}

\usepackage[T1]{fontenc}
\usepackage[utf8]{inputenc}
\DeclareFontShape{T1}{cmss}{b}{n}{<->ssub*cmss/bx/n}{}
\usepackage{graphicx}
\usepackage{amsmath,amssymb,bm,mathtools}
\usepackage{booktabs,tabularx,array}
\usepackage{xcolor}
\usepackage{tcolorbox}
\tcbuselibrary{skins}
\usepackage{hyperref}

% These prefixes allow compilation from either the repository root or this folder.
\graphicspath{{./}{../}{}}
\hypersetup{colorlinks=true,urlcolor=blue!55!black}

\newcommand{\hi}[1]{\textcolor{blue!55!black}{\textbf{#1}}}
\newcommand{\warn}[1]{\textcolor{red!70!black}{\textbf{#1}}}
\newcommand{\soft}[1]{\textcolor{gray!50!black}{#1}}
\newcommand{\xvec}{\bm{x}}
\newcommand{\Xvec}{\bm{X}}
\newcommand{\uvec}{\bm{u}}
\newcommand{\fvec}{\bm{f}}
\newcommand{\Ften}{\bm{F}}
\newcommand{\Pten}{\bm{P}^{e}}
\newcommand{\stress}{\bm{\sigma}}
\newcommand{\estress}{\bm{\sigma}^{e}}
\newcommand{\lforce}{\bm{\mathcal{F}}}
\newcommand{\normal}{\bm{N}}
\newcommand{\currentnormal}{\bm{n}}
\newcommand{\Identity}{\bm{I}}

\newtcolorbox{punchbox}{
  enhanced,
  colback=blue!4,
  colframe=blue!55!black,
  boxrule=0.6pt,
  arc=2pt,
  left=4pt,right=4pt,top=3pt,bottom=3pt
}

\title[IBFE motivation]{Why combine immersed boundary and finite element methods?}
\subtitle{The mechanics, coupling, and tradeoffs of IBFE}
\author[Miller]{Laura A. Miller}
\institute[UArizona]{Department of Mathematics, University of Arizona}
\date{}

\begin{document}

% ----- 1. Title -----
\begin{frame}[plain]
  \begin{beamercolorbox}[sep=10pt,center]{title}
    {\bfseries\Large Why combine immersed boundary and finite element methods?}\\[3pt]
    {\normalsize The mechanics, coupling, and tradeoffs of IBFE}
  \end{beamercolorbox}
  \vspace{0.8em}
  \begin{columns}[c,onlytextwidth]
  \begin{column}{0.38\textwidth}
    \textbf{\large Laura A. Miller}\\[4pt]
    Department of Mathematics\\
    University of Arizona\\[13pt]
    {\small IBFE Tutorials}
  \end{column}
  \begin{column}{0.62\textwidth}
    \centering
    \includegraphics[width=0.95\linewidth]{pictures/4-3/IBFE-sketch.png}
  \end{column}
  \end{columns}
  \vfill
  \centering{\tiny\color{black!55}AI tools were used to assist with formatting, organization, and clarification of the instructor's original notes.}
\end{frame}

% ----- 2. Roadmap -----
\begin{frame}{Roadmap}
\large
\begin{enumerate}\itemsep10pt
  \item Why a fixed fluid grid helps with large structural motion
  \item How finite elements represent the solid mechanics
  \item How integral operators exchange velocity and force
  \item Which approximations and tradeoffs matter in practice
\end{enumerate}

\vfill
\begin{punchbox}
\centering
IBFE keeps the immersed interface treatment and uses a variational finite element model for solid mechanics.
\end{punchbox}
\end{frame}

% ----- 3. FSI challenge -----
\begin{frame}{The fluid--structure interaction challenge}
\begin{columns}[T,onlytextwidth]
\begin{column}{0.48\textwidth}
\hi{Fluid description}
\begin{itemize}\itemsep6pt
  \item A fixed Eulerian grid supports efficient incompressible-flow solvers.
  \item Velocity and pressure live at spatial locations.
  \item The grid should remain well shaped as the structure moves.
\end{itemize}
\end{column}
\begin{column}{0.48\textwidth}
\hi{Solid description}
\begin{itemize}\itemsep6pt
  \item A material mesh follows deformation and rotation.
  \item Constitutive laws depend on strain relative to a reference state.
  \item Large motion can severely distort a body-fitted fluid mesh.
\end{itemize}
\end{column}
\end{columns}

\medskip
\begin{punchbox}
The numerical method must enforce no slip and transmit elastic force without requiring the two meshes to conform.
\end{punchbox}
\end{frame}

% ----- 4. Configurations: retain the original figure occurrence -----
\begin{frame}{Reference and current configurations}
\begin{columns}[c,onlytextwidth]
\begin{column}{0.59\textwidth}
  \centering
  \includegraphics[width=0.98\linewidth]{pictures/lagrangian-eulerian-coordinate.png}

  {\scriptsize Reference configuration \(B_0\), current configuration \(B_t\), and the motion \(\bm{\chi}\). Figure 1 of Griffith \& Luo (2017), \emph{Int. J. Numer. Methods Biomed. Eng.} 33:e2888. NB: the figure labels the reference domain \(U\); in this course it is \(B_0\).}
\end{column}
\begin{column}{0.39\textwidth}
  \[
    \bm{\chi}:B_0\times[0,T]\longrightarrow\Omega
  \]
  \[
    \xvec=\bm{\chi}(\Xvec,t)
  \]
  \begin{itemize}\itemsep3pt
    \item \(B_0\): reference solid domain
    \item \(B_t=\bm{\chi}(B_0,t)\): current solid region
    \item \(\Omega\): fixed Eulerian domain
    \item \(\Xvec\): material coordinate
    \item \(\xvec\): spatial coordinate
  \end{itemize}
\end{column}
\end{columns}

\smallskip
\begin{punchbox}
\scriptsize
Notation change from Module 1: there the current position of material point \(q\) was written \(\bm{X}(q,t)\). From here on we follow Griffith \& Luo: \(\bm{X}\) is the reference point in \(B_0\) and \(\bm{\chi}(\bm{X},t)\) is its current position in \(B_t\); \(\Omega\) is the Eulerian computational domain.
\end{punchbox}
\end{frame}

% ----- 5. Method landscape -----
\begin{frame}{Interface treatment in common FSI methods}
\small
\renewcommand{\arraystretch}{1.35}
\begin{tabularx}{\textwidth}{
  >{\bfseries\raggedright\arraybackslash}p{0.20\textwidth}
  >{\raggedright\arraybackslash}X
  >{\raggedright\arraybackslash}X}
\toprule
Method & Fluid mesh & Structural mechanics \\
\midrule
Body-fitted / ALE
  & Conforms to the interface. Mesh motion or remeshing may be needed
  & Standard finite elements and sharp boundary conditions \\
Classical IB
  & Fixed Cartesian grid with regularized interface coupling
  & Lagrangian fibers, springs, or continuum energy models \\
IBFE
  & Fixed Cartesian grid with regularized interface coupling
  & Standard finite elements with weak-form stress divergence \\
\bottomrule
\end{tabularx}

\medskip
\begin{punchbox}
IBFE couples sophisticated solid mechanics to a fixed, nonconforming fluid grid while retaining a regularized interface.
\end{punchbox}
\end{frame}

% ----- 6. Hybrid discretization -----
\begin{frame}{The hybrid discretization}
\vspace{0.4em}
\begin{beamercolorbox}[wd=\textwidth,sep=9pt,center,rounded=true]{block body}
{\Large \hi{Eulerian fluid solver}}
\qquad \(+\) \qquad
{\Large \hi{Lagrangian FE solid}}
\end{beamercolorbox}

\vspace{1.0em}
\begin{columns}[T,onlytextwidth]
\begin{column}{0.47\textwidth}
\textbf{Finite differences or finite volumes}
\begin{itemize}\itemsep5pt
  \item incompressible momentum equation
  \item pressure projection or saddle-point solve
  \item Cartesian-grid adaptivity when needed
\end{itemize}
\end{column}
\begin{column}{0.47\textwidth}
\textbf{Finite elements}
\begin{itemize}\itemsep5pt
  \item large-deformation kinematics
  \item anisotropic or nonlinear material response
  \item consistent weak-form elastic forces
\end{itemize}
\end{column}
\end{columns}

\vfill
\centering
\hi{Regularized delta kernels connect the two discretizations.}
\end{frame}

% ----- 7. One-field continuum model -----
\begin{frame}{One velocity field for the coupled system}
\begin{columns}[c,onlytextwidth]
\begin{column}{0.51\textwidth}
The original IBFE formulation uses a composite Eulerian velocity field \(\uvec(\xvec,t)\) throughout \(\Omega\).

\medskip
The total Cauchy stress is
\[
\stress(\xvec,t)=
\begin{cases}
\stress^f(\xvec,t)+\estress(\xvec,t), & \xvec\in B_t,\\
\stress^f(\xvec,t), & \xvec\notin B_t.
\end{cases}
\]
\end{column}
\begin{column}{0.45\textwidth}
\begin{punchbox}
The background viscous stress acts everywhere. The immersed solid contributes an additional elastic stress inside \(B_t\).
\end{punchbox}

\medskip
For a Newtonian background fluid,
\[
\stress^f=-p\Identity
+\mu\left(\nabla\uvec+\nabla\uvec^{T}\right).
\]
\end{column}
\end{columns}

\vfill
{\scriptsize Scope of Griffith--Luo (2017): common density and background viscosity. Later extensions treat material contrasts.}
\end{frame}

% ----- 8. Governing system -----
\begin{frame}{Continuum equations}
\begin{columns}[c,onlytextwidth]
\begin{column}{0.56\textwidth}
\begin{align*}
\rho\left(\frac{\partial\uvec}{\partial t}
+\uvec\!\cdot\!\nabla\uvec\right)
  &= -\nabla p+\mu\nabla^2\uvec+\fvec^{e},\\[2pt]
\nabla\!\cdot\uvec&=0,\\[2pt]
\frac{\partial\bm{\chi}}{\partial t}(\Xvec,t)
  &=\uvec(\bm{\chi}(\Xvec,t),t).
\end{align*}
\end{column}
\begin{column}{0.40\textwidth}
\textbf{Three coupled ingredients}
\begin{itemize}\itemsep7pt
  \item incompressible momentum balance
  \item Eulerian elastic force \(\fvec^e\)
  \item material motion with the local fluid velocity
\end{itemize}
\end{column}
\end{columns}

\medskip
\begin{punchbox}
At the continuum level, the last equation imposes both no slip and no penetration on the immersed material.
\end{punchbox}
\end{frame}

% ----- 9. Kinematic coupling -----
\begin{frame}{Kinematic coupling as an integral transform}
\[
\frac{\partial\bm{\chi}}{\partial t}(\Xvec,t)
=\int_{\Omega}\uvec(\xvec,t)\,
\delta\!\left(\xvec-\bm{\chi}(\Xvec,t)\right)\,d\xvec.
\]

\begin{columns}[T,onlytextwidth]
\begin{column}{0.47\textwidth}
\hi{Continuum identity}

The Dirac distribution evaluates the Eulerian velocity at the current material position.
\end{column}
\begin{column}{0.47\textwidth}
\hi{Discrete approximation}

A regularized kernel \(\delta_h\) averages velocity over several grid cells. No slip is therefore enforced in a regularized sense.
\end{column}
\end{columns}

\vfill
\begin{punchbox}
The kernel width controls the numerical interface thickness. It does not represent a physical coating around the solid.
\end{punchbox}
\end{frame}

% ----- 10. Stress measures -----
\begin{frame}{Stress and traction in two configurations}
\begin{columns}[T,onlytextwidth]
\begin{column}{0.48\textwidth}
\hi{Current configuration}
\[
\bm{t}=\stress\currentnormal
\]
The Cauchy traction is force per unit current area. Its normal and tangential components describe pressure-like and shear loading.
\end{column}
\begin{column}{0.48\textwidth}
\hi{Reference configuration}
\[
\bm{t}_0=\Pten\normal
\]
The first Piola--Kirchhoff traction \(\bm{t}_0\) is force per unit reference area. It is the natural surface quantity for a Lagrangian FE weak form.

\smallskip
{\scriptsize (we write \(\bm{P}^{e}\) to stress that it is the elastic stress; later decks drop the superscript)}
\end{column}
\end{columns}

\medskip
\begin{punchbox}
The fluid traction on a physical interface and the elastic transmission force used in the IB equations are related by momentum balance, but they are not interchangeable labels.
\end{punchbox}
\end{frame}

% ----- 11. Piola transformation -----
\begin{frame}{Why IBFE uses the first Piola--Kirchhoff stress}
\begin{columns}[c,onlytextwidth]
\begin{column}{0.55\textwidth}
\[
\Ften=\nabla_{\!X}\bm{\chi},
\qquad J=\det\Ften,
\]
\[
\boxed{\Pten(\Xvec,t)
=J\,\estress(\bm{\chi}(\Xvec,t),t)\,\Ften^{-T}}
\]
\end{column}
\begin{column}{0.41\textwidth}
\begin{itemize}\itemsep7pt
  \item \(\Ften\) maps material line elements to the current configuration.
  \item \(J\) gives the local volume ratio.
  \item \(\Pten\) need not be symmetric because it links reference normals to current forces.
\end{itemize}
\end{column}
\end{columns}

\medskip
\begin{punchbox}
The Piola transformation makes the reference-domain and current-domain traction resultants agree.
\end{punchbox}
\end{frame}

% ----- 12. Strong force -----
\begin{frame}{Elastic stress becomes an Eulerian force}
\small
\[
\boxed{
\begin{aligned}
\fvec^{e}(\xvec,t)
={}&\int_{B_0}\operatorname{Div}_{X}\Pten(\Xvec,t)\,
\delta\!\left(\xvec-\bm{\chi}(\Xvec,t)\right)\,d\Xvec\\[-1pt]
&-\int_{\partial B_0}\Pten(\Xvec,t)\normal\,
\delta\!\left(\xvec-\bm{\chi}(\Xvec,t)\right)\,dA.
\end{aligned}}
\]

\begin{columns}[T,onlytextwidth]
\begin{column}{0.47\textwidth}
\hi{Internal force}

\(\operatorname{Div}_{X}\Pten\) acts throughout the reference solid and has units of force per reference volume.
\end{column}
\begin{column}{0.47\textwidth}
\hi{Transmission force}

\(-\Pten\normal\) acts on \(\partial B_0\) and has units of force per reference area. It produces interfacial stress jumps.
\end{column}
\end{columns}

\vfill
{\scriptsize Source: Griffith and Luo (2017), Eq.~(6) for the boxed force equation (strong form, Eqs. (4)--(7)).}
\end{frame}

% ----- 13. Unified weak form -----
\begin{frame}{The unified weak force}
Rather than differentiate a piecewise FE stress, define a Lagrangian force density \(\lforce\) by
\[
\boxed{
\int_{B_0}\lforce\!\cdot\!\bm{V}\,d\Xvec
=-\int_{B_0}\Pten:\nabla_{\!X}\bm{V}\,d\Xvec
\qquad\text{for every FE test function }\bm{V}.}
\]
Then spread the result to the fluid:
\[
\fvec^{e}(\xvec,t)=\int_{B_0}\lforce(\Xvec,t)\,
\delta\!\left(\xvec-\bm{\chi}(\Xvec,t)\right)\,d\Xvec.
\]

\begin{punchbox}
The test function does not vanish on \(\partial B_0\). Consequently, this one weak equation contains both the internal and boundary transmission forces.
\end{punchbox}
\end{frame}

% ----- 14. Two weak formulations -----
\begin{frame}{Unified and partitioned weak formulations}
\begin{columns}[T,onlytextwidth]
\begin{column}{0.48\textwidth}
\hi{Unified form}
\begin{itemize}\itemsep6pt
  \item One volumetric weak force represents internal and boundary terms.
  \item Adjoint coupling gives an exact discrete interaction-power identity.
  \item The boundary contribution enters through the volume representation.
\end{itemize}
\end{column}
\begin{column}{0.48\textwidth}
\hi{Partitioned form}
\begin{itemize}\itemsep6pt
  \item Internal force and \(\bm{T}=-\Pten\normal\) are spread separately (the reference traction is \(\bm{t}_0=\Pten\normal\); the transmission force is its negative).
  \item Coarse solid meshes can give better pressure accuracy and volume conservation.
  \item The same discrete power identity does not hold.
\end{itemize}
\end{column}
\end{columns}

\medskip
\begin{punchbox}
The two forms are equivalent before discretization, but their numerical solutions need not coincide.
\end{punchbox}

\vfill
{\scriptsize Source: Griffith and Luo (2017), Sections 2.2 and 7.}
\end{frame}

% ----- 15. FE approximation -----
\begin{frame}{Finite elements carry the solid kinematics}
\begin{columns}[c,onlytextwidth]
\begin{column}{0.56\textwidth}
\begin{align*}
\bm{\chi}_h(\Xvec,t)
  &=\sum_{a=1}^{M}\bm{\chi}_a(t)\,\phi_a(\Xvec),\\[3pt]
\Ften_h(\Xvec,t)
  &=\nabla_{\!X}\bm{\chi}_h
   =\sum_{a=1}^{M}\bm{\chi}_a(t)\otimes\nabla_{\!X}\phi_a.
\end{align*}
\end{column}
\begin{column}{0.40\textwidth}
\begin{itemize}\itemsep7pt
  \item Nodal positions define a continuous configuration.
  \item \(C^0\) elements generally give elementwise-discontinuous \(\Ften_h\) and \(\Pten_h\).
  \item The weak form avoids differentiating \(\Pten_h\) across element boundaries.
\end{itemize}
\end{column}
\end{columns}

\medskip
For basis function \(\phi_b\), the unified force load is
\[
B_b=-\int_{B_0}\Pten_h:\nabla_{\!X}\phi_b\,d\Xvec,
\qquad \bm{M}[\lforce]=[\bm{B}].
\]
\end{frame}

% ----- 16. Quadrature coupling -----
\begin{frame}{Quadrature points mediate force transfer}
\small
For each finite element \(B_0^e\), evaluate the FE force density at quadrature points \(\Xvec_Q^e\):
\[
\boxed{
\fvec_h(\xvec,t)=
\sum_{B_0^e}\sum_{Q=1}^{N_e}
\lforce_h(\Xvec_Q^e,t)\,
\delta_h\!\left(\xvec-\bm{\chi}_h(\Xvec_Q^e,t)\right)\,w_Q^e.}
\]

\begin{columns}[T,onlytextwidth]
\begin{column}{0.47\textwidth}
\hi{Classical nodal coupling}

Sparse nodes can leave gaps in kernel support and permit numerical leakage.
\end{column}
\begin{column}{0.47\textwidth}
\hi{IBFE quadrature coupling}

The element mesh may be coarser than the fluid grid if the quadrature-point net remains sufficiently dense.
\end{column}
\end{columns}

\vfill
{\scriptsize Structural resolution still depends on the FE mesh. Denser quadrature improves coupling resolution, not the constitutive approximation.}
\end{frame}

% ----- 17. Adjoint coupling -----
\begin{frame}{Adjoint coupling preserves interaction power}
Write the discrete force-spreading and velocity-restriction operations as
\[
\fvec_h=\mathcal{S}[\bm{\chi}_h]\lforce_h,
\qquad
\frac{d\bm{\chi}_h}{dt}=\mathcal{R}[\bm{\chi}_h]\uvec_h.
\]

Choose \(\mathcal{R}=\mathcal{S}^{*}\) with respect to the Eulerian and FE inner products. Then
\[
\boxed{
\left(\lforce_h,\frac{d\bm{\chi}_h}{dt}\right)_{B_0}
=\left(\fvec_h,\uvec_h\right)_{\Omega}.}
\]

\begin{punchbox}
The Eulerian and Lagrangian descriptions assign the same instantaneous power to the coupling force. This property alone does not guarantee stability of the full time integrator.
\end{punchbox}
\end{frame}

% ----- 18. Computational cycle -----
\begin{frame}{One IBFE time step}
\begin{enumerate}\itemsep7pt
  \item Restrict the Eulerian velocity to FE degrees of freedom.
  \item Update the current configuration \(\bm{\chi}_h\).
  \item Evaluate \(\Ften_h\), \(J_h\), and the constitutive stress \(\Pten_h\).
  \item Assemble the Lagrangian weak force coefficients.
  \item Spread quadrature-point forces to the Cartesian grid and advance the fluid.
\end{enumerate}

\medskip
\begin{punchbox}
Partitioned, semi-implicit, and fully implicit algorithms reorder or couple these operations differently. The transfer operators and variational solid mechanics remain the organizing ideas.
\end{punchbox}
\end{frame}

% ----- 19. Material models -----
\begin{frame}{Constitutive models in the FE solid}
\begin{columns}[T,onlytextwidth]
\begin{column}{0.49\textwidth}
\hi{Hyperelastic example}
\[
\Pten_{\mathrm{pas}}=\frac{\partial W}{\partial\Ften}
\]
The strain-energy density \(W\) may represent isotropic, fiber-reinforced, or nearly incompressible response.
\end{column}
\begin{column}{0.47\textwidth}
\hi{More general stress laws}

The numerical framework needs a routine that evaluates a Lagrangian stress. It can also accommodate history-dependent or viscoelastic response when the required internal variables are evolved.
\end{column}
\end{columns}

\medskip
\begin{punchbox}
Griffith and Luo use hyperelastic examples for simplicity, but the method itself does not require an explicit strain-energy functional.
\end{punchbox}
\end{frame}

% ----- 20. Active stress -----
\begin{frame}{Active tension belongs in the stress model}
For a passive material with an active fiber response, decompose
\[
\Pten=\Pten_{\mathrm{pas}}+\Pten_{\mathrm{act}}.
\]

\begin{columns}[T,onlytextwidth]
\begin{column}{0.48\textwidth}
\textbf{A common construction}
\begin{itemize}\itemsep6pt
  \item \(\bm{a}=\Ften\bm{f}_0\) (mapped fiber; \(\bm{f}_0\) a unit reference direction). Note \(\bm{a}\) is \emph{not} normalized.
  \item prescribe an active Cauchy stress:
  \(\estress_{\mathrm{act}}=T\,\bm{a}\otimes\bm{a}\).
  \item map it to the reference configuration:
  \(\Pten_{\mathrm{act}}=J\estress_{\mathrm{act}}\Ften^{-T}
   =J\,T\,\bm{a}\otimes\bm{f}_0\).
\end{itemize}
\end{column}
\begin{column}{0.48\textwidth}
\textbf{Model choices}
\begin{itemize}\itemsep6pt
  \item \(T\) may depend on activation, stretch, and stretch rate.
  \item Other conventions normalize \(\bm{a}\) or absorb stretch factors into \(T\); see deck 4-3.
\end{itemize}
\end{column}
\end{columns}

\medskip
\begin{punchbox}
A stress-based activation law makes material-frame invariance and local momentum balance explicit.
\end{punchbox}
\end{frame}

% ----- 21. Volume conservation -----
\begin{frame}{Incompressibility and structural volume}
At the continuum level, material motion with a divergence-free velocity gives
\[
\frac{dJ}{dt}=J\,\nabla\!\cdot\uvec(\bm{\chi}(\Xvec,t),t)=0.
\]

\begin{columns}[T,onlytextwidth]
\begin{column}{0.47\textwidth}
\hi{Continuum consequence}

An initially volume-preserving deformation remains volume preserving under the exact incompressible velocity field.
\end{column}
\begin{column}{0.47\textwidth}
\hi{Discrete reality}

Kernel regularization, velocity projection, time integration, and under-resolution can introduce volume drift.
\end{column}
\end{columns}

\medskip
\begin{punchbox}
Mesh refinement, compatible coupling, appropriate quadrature, and volumetric stabilization can reduce drift. None replaces a convergence study.
\end{punchbox}
\end{frame}

% ----- 22. Strengths -----
\begin{frame}{What IBFE does especially well}
\begin{itemize}\itemsep8pt
  \item Handles large translations, rotations, and deformations without remeshing the fluid grid.
  \item Reuses standard finite element ideas for nonlinear and anisotropic solid mechanics.
  \item Separates structural mesh spacing from coupling-point density through element quadrature.
  \item Fits naturally with Cartesian-grid solvers and block-structured adaptive mesh refinement.
\end{itemize}

\vfill
\begin{punchbox}
The main benefit is a flexible division of labor: the Eulerian discretization handles incompressible flow, while the FE discretization handles material mechanics.
\end{punchbox}
\end{frame}

% ----- 23. Limitations -----
\begin{frame}{Persistent numerical challenges}
\begin{itemize}\itemsep7pt
  \item The regularized kernel spreads the interface over several grid cells, which limits sharp stress and pressure resolution.
  \item Stiff elastic response can impose severe time-step restrictions in explicit or loosely coupled schemes.
  \item Contact, self-contact, tearing, and topology change require additional models or algorithms.
  \item Density and viscosity contrast require extensions beyond the simplest one-field formulation.
  \item Accurate results still require adequate fluid resolution, FE resolution, and coupling quadrature.
\end{itemize}

\medskip
\begin{punchbox}
IBFE improves geometric flexibility, but sharp interfaces and stiff material laws remain costly.
\end{punchbox}
\end{frame}

% ----- 24. Method selection -----
\begin{frame}{When IBFE is a good match}
\begin{columns}[T,onlytextwidth]
\begin{column}{0.48\textwidth}
\hi{Strong fit}
\begin{itemize}\itemsep6pt
  \item large solid motion or deformation
  \item detailed nonlinear constitutive behavior
  \item fixed material topology
  \item tolerance for a regularized interface
\end{itemize}
\end{column}
\begin{column}{0.48\textwidth}
\hi{Consider another formulation when}
\begin{itemize}\itemsep6pt
  \item pointwise surface stresses or thin boundary layers dominate the objective
  \item rupture or changing topology is central
  \item a reliable body-fitted mesh remains inexpensive to maintain
\end{itemize}
\end{column}
\end{columns}

\medskip
\begin{punchbox}
Method choice should follow the required observables: bulk motion and deformation, or sharply resolved interfacial fields.
\end{punchbox}
\end{frame}

% ----- 25. Applications -----
\begin{frame}{Representative IBFE problems}
\begin{columns}[T,onlytextwidth]
\begin{column}{0.48\textwidth}
\textbf{Benchmarks in Griffith--Luo}
\begin{itemize}\itemsep6pt
  \item anisotropic elastic shells
  \item an elastic disc in a lid-driven cavity
  \item pressure-driven collapsible-channel flow
  \item the Turek--Hron cylinder-and-flap benchmark
\end{itemize}
\end{column}
\begin{column}{0.48\textwidth}
\textbf{Broader application areas}
\begin{itemize}\itemsep6pt
  \item heart valves and cardiovascular flow
  \item deformable biological tissues
  \item flexible structures in viscous flow
  \item actively contracting materials
\end{itemize}
\end{column}
\end{columns}

\vfill
{\scriptsize See Griffith and Luo (2017) and Griffith and Patankar, \emph{Annu. Rev. Fluid Mech.} 2020.}
\end{frame}

% ----- 26. Resources -----
\begin{frame}{Resources}
\begin{itemize}\itemsep8pt
  \item \href{https://doi.org/10.1002/cnm.2888}{Griffith and Luo (2017): hybrid finite difference / finite element IB method}
  \item \href{https://doi.org/10.1146/annurev-fluid-010719-060228}{Griffith and Patankar (2020): immersed methods for FSI}
  \item \href{https://ibamr.github.io/}{IBAMR software and documentation}
  \item \href{https://www.youtube.com/watch?v=-1TlNJL93sc}{Boyce Griffith presentation on IBFE}
  \item \href{https://drive.google.com/drive/folders/1vkhrX0YInMm5pfUj3LcRA9TV0DZdFp1a?usp=sharing}{Solid mechanics and IBFE tutorial by Matea}
  \item \href{https://www.classcentral.com/classroom/youtube-aim-micde-seminar-boyce-griffith-immersed-methods-for-fluid-structure-interaction-323154}{AIM/MICDE seminar: immersed methods for FSI}
\end{itemize}
\end{frame}

% ----- 27. Closing -----
\begin{frame}{Three ideas to carry forward}
\begin{enumerate}\itemsep11pt
  \item \hi{Mechanics:} finite elements provide the deformation map, stress, and weak elastic force.
  \item \hi{Coupling:} regularized integral operators exchange force and velocity between nonconforming grids.
  \item \hi{Numerics:} quadrature density, adjointness, and resolution determine whether the hybrid method realizes its continuum properties.
\end{enumerate}

\vfill
\begin{beamercolorbox}[wd=\textwidth,sep=10pt,center,rounded=true]{title}
{\Large Questions?}
\end{beamercolorbox}
\end{frame}

\end{document}
